%; whizzy chapter
% -initex iniptex -latex platex -format platex -bibtex jbibtex -fmt fmt
% 以上 whizzytex を使用する場合の設定。

%     Kansai Debian Meeting resources
%     Copyright (C) 2007 Takaya Yamashita
%     Thank you for Tokyo Debian Meeting resources

%     This program is free software; you can redistribute it and/or modify
%     it under the terms of the GNU General Public License as published by
%     the Free Software Foundation; either version 2 of the License, or
%     (at your option) any later version.

%     This program is distributed in the hope that it will be useful,
%     but WITHOUT ANY WARRANTY; without even the implied warranty of
%     MERCHANTABILITY or FITNESS FOR A PARTICULAR PURPOSE.  See the
%     GNU General Public License for more details.

%     You should have received a copy of the GNU General Public License
%     along with this program; if not, write to the Free Software
%     Foundation, Inc., 51 Franklin St, Fifth Floor, Boston, MA  02110-1301 USA

%  preview (shell-command (concat "evince " (replace-regexp-in-string "tex$" "pdf"(buffer-file-name)) "&"))
% 画像ファイルを処理するためにはebbを利用してboundingboxを作成。
%(shell-command "cd image200708; ebb *.png")

%%ここからヘッダ開始。

\documentclass[mingoth,a4paper]{jsarticle}
\usepackage{kansaimonthlyreport}
\usepackage[dvipdfmx]{xy}
\usepackage{etex}
\usepackage{ulem}

% 日付を定義する、毎月変わります。
\newcommand{\debmtgyear}{2016}
\newcommand{\debmtgdate}{27}
\newcommand{\debmtgmonth}{3}
\newcommand{\debmtgnumber}{108}

\def\fixme#1{{\color{red}{#1}}}

\begin{document}

\begin{titlepage}

% 毎月変更する部分、本文の末尾も修正することをわすれずに

  第\debmtgnumber{}回 関西 Debian 勉強会資料

  \vspace{2cm}

  \begin{center}
    \includegraphics{image200802/kansaidebianlogo.png}
  \end{center}

  \begin{flushright}
    \hfill{}関西Debian勉強会担当者 佐々木・倉敷・のがた・かわだ \\
    \hfill{}\debmtgyear{}年\debmtgmonth{}月\debmtgdate{}日
  \end{flushright}

  \thispagestyle{empty}
\end{titlepage}

\dancersection{Introduction}{Debian JP}

\vspace{1em}

関西Debian勉強会はDebian GNU/Linuxのさまざまなトピック
(新しいパッケージ、Debian特有の機能の仕組、Debian界隈で起こった出来事、
などなど）について話し合う会です。

目的として次の三つを考えています。
\begin{itemize}
\item MLや掲示板ではなく、直接顔を合わせる事での情報交換の促進
\item 定期的に集まれる場所
\item 資料の作成
\end{itemize}

 それでは、楽しい一時をお過ごしください。

\newpage

\begin{minipage}[b]{0.2\hsize}
  {\rotatebox{90}{\fontsize{80}{80}
{\gt{関西 Debian 勉強会}}}}
\end{minipage}
\begin{minipage}[b]{0.8\hsize}
\hrule
\vspace{2mm}
\hrule
\setcounter{tocdepth}{1}
\tableofcontents
\vspace{2mm}
\hrule
\end{minipage}

\dancersection{最近のDebian関係のイベント報告}{Debian JP}

\subsection{第107回関西Debian勉強会}

107回目の関西Debian勉強会は2月28日(日)に福島区民センターで開催されました。

かわださんによる「勉強会資料の歩き方」と倉敷さんによる「周回遅れでDocker触ってみ
た」の二本立てでした。

\subsection{第138回東京エリアDebian勉強会}

138回目の東京エリアDebian勉強会は3月5日(土)にサイボウズ株式会社東京オフィスで開
催されました。

内容はAdrianさんによる「How to become a Debian Developer」とwskokaさんによる
「tilegxについて(仮)」、kenhysさんによる「Debianのインフラを借りるには」の三本立
ての内容でした。

\subsection{Debian Project}

\subsubsection {Renaming Iceweasel to Firefox}

みなさんご存知のとおりIceweaselがFirefoxに戻りました。

Sylvestre Ledruさんが\debianbug{815006}に登録したのはつぎのような内容でした。

\begin{itemize}
\item MozillaはDebian固有のパッチがIceweasel/Firefoxの品質に影響ないと認識する
\item Debian stableリリース中でもFirefox ESRのバージョンは更新する
\item \debianbug{354622}の問題は解消され、FirefoxのロゴはDFSGに合致する
\item jessieではメンテナンスのためIceweaselの名は残る
\item Icedove/Thunderbirdについても同様の対応をすすめる
\end{itemize}

そして2016-03-11にfirefox-esr 45.0esr-1とfirefox 45.0-1パッケージがunstableに入
りました。firefox 1.5.dfsg+1.5.0.7-2がtestingが取り除かれたのが2007-01-18ですか
ら9年2ヶ月ぶりにDebianに戻ってきたことになります。

また先のバグの内容からすると、stableにはfirefoxパッケージは入らず、firefox-esr
パッケージのみが入るのではないかと思われます。

\subsubsection {DebConf17 to be held in Montreal (CA)}

DebConfチームより来年のDebConf17の開催地がカナダのモントリオールに決定したとアナ
ウンスがありました%
\footnote{\url{https://lists.debian.org/debian-devel-announce/2016/02/msg00005.html}}。

カナダでDebConfが開催されるのはトロントで開催されたDebConf2以来15年ぶり2回目とな
ります。

\subsubsection{Bits from the Release Team: A Slightly Moveable Feast}

リリースチームよりフリーズの日程が少しずれるとアナウンスがありました%
\footnote{\url{https://lists.debian.org/debian-devel-announce/2016/03/msg00000.html}}。

Linux Kernelの長期サポートがみこまれる4.10をStretchに採用するためです。

全ての予定は二ヶ月遅れて次のとおりとなります。

\begin{itemize}
\item 2016/11/05 transitions freeze
\item 2017/01/05 "Soft" freeze
\item 2017/02/05 Full freeze
\end{itemize}

\subsubsection {Archive changes}

FTP MasterのJoerg Jaspertさんからarchiveの変更についてアナウンスがありました%
\footnote{\url{https://lists.debian.org/debian-devel-announce/2016/03/msg00006.html}}。

InRelease/Releaseファイル、Packages、SourcesからMD5とSHA1チェックサムがなくなり
SHA256チェックサムのみになり、gzipではなくxzでの提供となります。

まずはexperimentalのみ変更し、問題なければ順次sid、testingを変更して次のリリース
でstableも変更される予定でした。

しかし、jigdoがMD5を参照しておりすぐには対応できない%
\footnote{\url{https://lists.debian.org/debian-devel/2016/03/msg00194.html}}、
sources.debian.netも壊れてしまう \debianbug{818324} などの反応があり、ひとまず
SHA1チェックサムをなくす変更だけがおこなわれました。


\subsubsection {Linux Kernel ABI report}

Ponomarenko AndreyさんからLinuxカーネルのABIトラッキングリポートを更新したよと投
稿がありました%
\footnote{\url{https://lists.debian.org/debian-devel/2016/03/msg00171.html}}。

FreeBSDやHurdなど他のカーネルのサポートも追加しないといった反応もありましたが、
glibcやGTK+、Qtなどのトラッキングリポートを作成されています。
\url{http://abi-laboratory.pro/tracker/}

\dancersection{事前課題}{Debian JP}

今回の課題は以下の通りです。
\begin{screen}
  \begin{enumerate}
  \item %
    systemdで起動するJessie環境を御用意下さい。
    幾つか設定を変更することがありますので、VM環境の方が無難かもしれません。

  \end{enumerate}
\end{screen}

参加者の皆さんの解答は以下の通りです:

\begin{prework}{ lurdan }
  \begin{enumerate}
  \item 16時くらいに参上いたします……
  \end{enumerate}
\end{prework}

\begin{prework}{ t3rkwd }
  \begin{enumerate}
  \item わかりました。sid環境になるとおもいます。
  \end{enumerate}
\end{prework}

\begin{prework}{ むんくさん }
  \begin{enumerate}
  \item 承知しました。Note PC 持参します。
  \end{enumerate}
\end{prework}

\begin{prework}{ 佐々木洋平 }
  \begin{enumerate}
  \item 無駄に、vagrant-kvm に挑戦してみたり。
  \end{enumerate}
\end{prework}

\begin{prework}{ tsukudamayo }
  \begin{enumerate}
  \item VMの環境です。よろしくお願いします。
  \end{enumerate}
\end{prework}

\begin{prework}{ 川江 浩 }
  \begin{enumerate}
  \item 了解。
  \end{enumerate}
\end{prework}

\begin{prework}{ Kazuhiro NISHIYAMA }
  \begin{enumerate}
  \item VirtualBox で環境を用意しました。
  \end{enumerate}
\end{prework}

\begin{prework}{ Say-no }
  \begin{enumerate}
  \item わかりました。
  \end{enumerate}
\end{prework}

\begin{prework}{ hiroshi morimoto }
  \begin{enumerate}
  \item よろしくお願いします。
  \end{enumerate}
\end{prework}

\begin{prework}{ Katsuki Kobayashi }
  \begin{enumerate}
  \item よろしくおねがいします。
  \end{enumerate}
\end{prework}

\dancersection{systemdに浸ってみた}{佐々木 洋平}

\begin{flushright}
  All Your Control Groups Are Belong To Us! - Lennart Poettering, Red Hat Inc. \\
  \url{https://www.youtube.com/watch?v=MSG4jW187Is}
\end{flushright}

\subsection{はじめに}

今回のお題は以下の通りです:
\begin{quotation}
  Jessie から default となった init である systemd ですが
  「全てがsystemd になる」と揶揄されるように他にも様々な機能があります。
  今回はパッケージで提供されている systemd 関連のツールを一通り試してみて、
  どう使うのか/使えるのかについての検証結果を報告します。
\end{quotation}

というわけで、
Jessie において提供されている幾つかの \verb|system-*| 関連のソフトウェアを使ってみます。
%
サービスの起動に関連した systemd の使い方については、
\cite[前回(2014年6月)の勉強会資料]{佐々木(2014)}をご覧下さい\footnote{もう二年前かよ$\dots$}。

\subsection{どんな環境のお話?}

調査対象は素の Debian ver.8 Jessie です。
「どのぐらい『素』か?」というとインストール時に「SSH サーバ」以外を外した状況です%
\footnote{%
  普段使いの環境は sid なので、%
  仮想環境として libvirt - KVM を利用しています。%
  LXC でも同様なのか、は良くわかりません。%
  親環境と子環境がともに systemd だった場合にも LXC はちゃんと動くようになったのかしら。%
  以前は祟りがあった様な$\dots$。どなたかご存知ですか?%
}。
\verb|dpkg -l| の出力結果は以下の通り
\begin{commandline}
  % dpkg -l | grep ^ii | wc -l
  279
\end{commandline}
\noindent%
まあ、minimal に比べれば多少多いですが。

さて、この段階で導入されているsystemd関連のパッケージはなんでしょうか。
\begin{commandline}
  % dpkg -l | grep systemd
   ii  libsystemd0:amd64  215-17+deb8u3  amd64  systemd utility library
   ii  systemd            215-17+deb8u3  amd64  system and service manager
   ii  systemd-sysv       215-17+deb8u3  amd64  system and service manager - SysV links
\end{commandline}
\noindent%
というわけで基本セットが入っている状況です。
しかしながら基本セットだけですと \verb|dbus| がインストールされていないため、
後述の幾つかの daemon が動きません%
\footnote{systemd のパッケージでは dbus は Recommends になっています}。
というわけで、\verb|dbus| と、旧来の sysvinit のスクリプトを上手く扱うために \verb|systemd-shim| をインストールしておきます。
\begin{commandline}
  % sudo apt-get install dbus systemd-shim --no-install-recommends
\end{commandline}

\subsection{時計合わせ - systemd-timesyncd}

時計合わせといえば NTP ですね。
systemd の流儀に従うのであれば \verb|systemd-timesyncd| を使うことになります。
とりあえず、現状を確認するために \verb|timedatectl| を叩いてみましょう%
\footnote{$\dots$ 執筆時刻がバレますね}。
\begin{commandline}
% timedatectl
      Local time: 日 2016-03-27 06:16:54 JST
  Universal time: 土 2016-03-26 21:16:54 UTC
        RTC time: 土 2016-03-26 21:16:54
       Time zone: Asia/Tokyo (JST, +0900)
     NTP enabled: yes
NTP synchronized: no
 RTC in local TZ: no
      DST active: n/a
\end{commandline}
\noindent
特に設定をしていませんが、\verb|NTP enabled: yes| です。
%
status を確認してみましょう。
\begin{commandline}
% systemctl status systemd-timesyncd
● systemd-timesyncd.service - Network Time Synchronization
   Loaded: loaded (/lib/systemd/system/systemd-timesyncd.service; enabled)
   Active: active (running) since 日 2016-03-27 06:24:11 JST; 2min 3s ago
     Docs: man:systemd-timesyncd.service(8)
 Main PID: 3255 (systemd-timesyn)
   Status: "Using Time Server 202.181.103.212:123 (0.debian.pool.ntp.org)."
   CGroup: /system.slice/systemd-timesyncd.service
           └─3255 /lib/systemd/systemd-timesyncd
\end{commandline}
\noindent
さて、\verb|Status: "Using Time Server 202.181.103.212:123 (0.debian.pool.ntp.org)."| とあるわけですが、これはどこから出てきた設定でしょうか?
%
実は、Debian の systemd パッケージはビルド時に
\begin{commandline}
  DEFAULT_NTP_SERVERS = 0.debian.pool.ntp.org 1.debian.pool.ntp.org 2.debian.pool.ntp.org 3.debian.pool.ntp.org
\end{commandline}
\noindent
を設定してビルドされています。
そんな訳で、特に設定をしなくても systemd-timesyncd は \verb|DEFAULT_NTP_SERVERS| を見に行くようになっています。

とはいえ、環境によって使いたい NTP サーバを変えたい事もあるでしょう。
そんな時は \verb|/etc/systemd/timesyncd.conf| の \verb|Servers| 行を修正します。
例えば
\begin{commandline}
#  This file is part of systemd.
- snip -

[Time]
Servers=ntp.kuins.kyoto-u.ac.jp
\end{commandline}
\noindent
として \verb|systemd-timesyncd| を再起動すると
\begin{commandline}
  % sudo systemctl daemon-reload
  % sudo systemctl restart systemd-timesyncd
  % systemctl status systemd-timesyncd
● systemd-timesyncd.service - Network Time Synchronization
   Loaded: loaded (/lib/systemd/system/systemd-timesyncd.service; enabled)
   Active: active (running) since 日 2016-03-27 06:33:38 JST; 1s ago
     Docs: man:systemd-timesyncd.service(8)
 Main PID: 3444 (systemd-timesyn)
   Status: "Using Time Server 130.54.240.26:123 (ntp.kuins.kyoto-u.ac.jp)."
   CGroup: /system.slice/systemd-timesyncd.service
           └─3444 /lib/systemd/systemd-timesyncd
  % timedatectl
      Local time: 日 2016-03-27 06:34:12 JST
  Universal time: 土 2016-03-26 21:34:12 UTC
        RTC time: 土 2016-03-26 21:34:11
       Time zone: Asia/Tokyo (JST, +0900)
     NTP enabled: yes
NTP synchronized: yes
 RTC in local TZ: no
      DST active: n/a
\end{commandline}
\noindent
といった塩梅になります。

時刻合わせの状況を知りたい時には少したってから status を叩いてみましょう。
\begin{commandline}
  % sudo systemctl status -l systemd-timesyncd
● systemd-timesyncd.service - Network Time Synchronization
   Loaded: loaded (/lib/systemd/system/systemd-timesyncd.service; enabled)
   Active: active (running) since 日 2016-03-27 08:37:02 JST; 5min ago
     Docs: man:systemd-timesyncd.service(8)
 Main PID: 1678 (systemd-timesyn)
   Status: ``Using Time Server 130.54.240.26:123 (ntp.kuins.kyoto-u.ac.jp).''
   CGroup: /system.slice/systemd-timesyncd.service
           └─1678 /lib/systemd/systemd-timesyncd

 3月 27 08:37:02 kvm-jessie-amd64 systemd-timesyncd[1678]: Using NTP server 130.54.240.26:123 (ntp.kuins.kyoto-u.ac.jp).
 3月 27 08:37:02 kvm-jessie-amd64 systemd-timesyncd[1678]: interval/delta/delay/jitter/drift 64s/-0.000s/0.002s/0.000s/+0ppm
 3月 27 08:38:06 kvm-jessie-amd64 systemd-timesyncd[1678]: interval/delta/delay/jitter/drift 128s/+0.001s/0.003s/0.000s/+1ppm
 3月 27 08:40:14 kvm-jessie-amd64 systemd-timesyncd[1678]: interval/delta/delay/jitter/drift 256s/-0.001s/0.003s/0.001s/+0ppm
\end{commandline}
\noindent
また、
複数のサーバを指定したい場合には \verb|Servers| 行に空白区切りでサーバを列挙します。

\subsubsection{別の NTP サービスを使っている場合は?}

ちょっと良くわかりません。
以前は \verb|/etc/systemd/ntp-units.d| 以下に
拡張子 \verb|.list| で使用しているサービス(Chrony なり NTPD なり) を記述していけば、
良きにはからってくれていたのですが$\dots$。

今回試した状況では、
\verb|systemd-timesyncd| と \verb|chrony| が両方起動して、
同時に別のサーバに時計を聞きに行く、といった挙動をしています。

非常に無駄な気もしますが、例えば、
\begin{itemize}
\item %
  chrony にローカルからの問い合わせに答える様にする
\item %
  timesyncd はローカルのchrony に問い合わせる
\end{itemize}
なんてことをすれば、まあ使えなくは無いです。

適切な NTP client が動作している場合には
そちらを見に行って \verb|NTP synchronized| を判断して欲しいのですが、
\sout{この辺の融通の効かなさが(ry}
\footnote{%
ちなみに、Ver.216 の changelog には
「もう /etc/systemd/ntp-units.d/ の下は見ねーからな? ちゃんと Conflicts 書けよ?」
って書いてあったりする。
}

\subsection{名前解決 - systemd-resolved}

続いて DNS を良きに図らう(?) systemd-resolved について。
%
とりえあず状況を確認すると$\dots$
\begin{commandline}
  % systemctl status systemd-resolved
● systemd-resolved.service - Network Name Resolution
   Loaded: loaded (/lib/systemd/system/systemd-resolved.service; disabled)
   Active: inactive (dead)
     Docs: man:systemd-resolved.service(8)
\end{commandline}
\noindent
というわけで、自動起動はされていません。
%
設定ファイルは /etc/systemd/resolved.conf で、ここに参照する DNS を空白区切りで記述します。
例えば
\begin{commandline}
  % cat /etc/systemd/resolved.conf
#  This file is part of systemd.
- snip -

[Resolve]
DNS=192.168.122.1
\end{commandline}
\noindent
として、おもむろに有効にしてみると$\dots$
\begin{commandline}
  % sudo systemctl enable systemd-resolved
  % sudo systemctl start systemd-resolved
\end{commandline}
\noindent
状況は以下の様に変化します:
\begin{commandline}
  % systemctl status systemd-resolved
● systemd-resolved.service - Network Name Resolution
   Loaded: loaded (/lib/systemd/system/systemd-resolved.service; enabled)
   Active: active (running) since 日 2016-03-27 09:13:34 JST; 1s ago
     Docs: man:systemd-resolved.service(8)
 Main PID: 1741 (systemd-resolve)
   Status: "Processing requests..."
   CGroup: /system.slice/systemd-resolved.service
           └─1741 /lib/systemd/systemd-resolved
\end{commandline}

この段階では /etc/resolv.conf の中身には変化がありません。
実の所、systemd-resolved の生成するファイルは\\
/run/systemd/resolve/resolv.conf にあるので、
/etc/resolv.conf をこちらへの symbolic link に変えておくと良いでしょう。
\begin{commandline}
  % cd /etc
  % sudo mv resolve.conf resolve.conf.bak
  % sudo sh -c 'ln -s /run/systemd/resolve/resolv.conf resolv.conf'
\end{commandline}
\noindent
これで名前解決ができるか確認しておきましょう。

ただ、現状では、これだけだと何が嬉しいのかさっぱりわかりません%
\footnote{%
  本来の目的は「caching and validationg DNS/DNSSEC stub resolver」なので%
  現状では機能が全然足りていない$\dots$
  sid にある systemd ver. 229 の man では御託宣がいろいろと増えているが、
  結局の所は動作に変更が無い様な。
}。
例えば
\begin{itemize}
\item %
  「caching」とあるのだけれど、cache の制御なんかは何処でやるの?
\item %
  \verb|search| 行が指定できないのだけれど、面倒じゃない?
\end{itemize}
とか。

一応、後述の \verb|systemd-networkd| と組み合わせるとネットワークに応じて DNS の切り替えができる、
というのは\textbf{設定ファイルが散らばらない}ので楽かもしれませんが%
\footnote{
$\dots$ しかし、オチがついた（笑）。systemd-networkd と組み合わせるのであれば、
DNS は networkd の設定ファイルで指定しておいた方が良いかも。
}。

\subsection{ネットワーク設定 - systemd-networkd}

お次はネットワーク設定を行なう systemd-networkd について。
とりあえず現状は
\begin{commandline}
  % systemctl status systemd-networkd
● systemd-networkd.service - Network Service
   Loaded: loaded (/lib/systemd/system/systemd-networkd.service; disabled)
   Active: inactive (dead)
     Docs: man:systemd-networkd.service(8)
\end{commandline}
\noindent
というわけで、これも自動起動はされていません。

多くの場合、有線ネットワークについては
/etc/network/interfaces に設定が書かれているのではないかと思うので、
その内容を元に systemd-networkd の設定ファイルを書いてみます。
設定ファイルとして、例えば
\begin{commandline}
  % cat /etc/systemd/network/wired.network
[Match]
Name=eth0

[Network]
Address=192.168.122.3/24
Gateway=192.168.122.1
DNS=192.168.122.1
# DNS=8.8.8.8     <-- DNS を複数指定する場合には、並べて書く。書かれた順に nameserver 行に並ぶ。
\end{commandline}
\noindent
なんてモノを用意します。
次に /etc/network/interfaces の中身を適当にコメントアウトして、
systemd-networkd を起動してみます。
\begin{commandline}
  % sudo systemctl enable systemd-networkd
  % sudo systemctl start systemd-networkd
  % /sbin/ifconfig
eth0      Link encap:イーサネット  ハードウェアアドレス 52:54:00:98:46:73
          inetアドレス:192.168.122.3 ブロードキャスト:192.168.122.255  マスク:255.255.255.0
          inet6アドレス: fe80::5054:ff:fe98:4673/64 範囲:リンク
          UP BROADCAST RUNNING MULTICAST  MTU:1500  メトリック:1
          RXパケット:14999 エラー:0 損失:0 オーバラン:0 フレーム:0
          TXパケット:10346 エラー:0 損失:0 オーバラン:0 キャリア:0
      衝突(Collisions):0 TXキュー長:1000
          RXバイト:4989650 (4.7 MiB)  TXバイト:1349343 (1.2 MiB)
...
\end{commandline}
という様になりました。
%
ついで /etc/resolv.conf の中身を確認してみると
\begin{commandline}
# This file is managed by systemd-resolved(8). Do not edit.
#
# Third party programs must not access this file directly, but
# only through the symlink at /etc/resolv.conf. To manage
# resolv.conf(5) in a different way, replace the symlink by a
# static file or a different symlink.

nameserver 192.168.122.1
nameserver 192.168.122.1
\end{commandline}
\noindent
$\dots$ えっと\textbf{重複は消してくれないんですかー?}

というわけで、/etc/systemd/resolved.conf の DNS 行 をコメントアウトしてから
おもむろに systemd-networkd を restart しましょう。

さて、これだけだとあまり嬉しくない(/etc/network/interfaces に書くのと変わらない)かもしれません。
もう少し設定を解説すると
\begin{itemize}
\item %
  Match では Name の他に「MACAddress」「Driver」「Host」なんてモノが使えます。
\item %
  Network では、当然「DHCP」「DHCPServer」等も使えます
\item %
  .network 以外にも .netdev, .link といった拡張子の設定ファイルが使えます。
  link の up/down や 仮想的な network device の追加/削除といった事が可能になります。
\end{itemize}
というわけで、
systemd が udev と連携していることから、
/etc/network/intefaces に書いていた時よりも、
より柔軟にネットワークの設定ができます(ことになっています)%
\footnote{%
  とはいえ、ちょっとまだ設定項目が足りない、かなぁ...。
  sid の systemd になると、もっと細かく制御できるのですが。
}。

\subsection{定期実行  - systemd-cron}

さて、大物です。systemd-cron はその名の通り cron の置き換えです。
おもむろに install しましょう
\begin{commandline}
  % sudo apt-get install systemd-cron --no-install-recommends
パッケージリストを読み込んでいます... 完了
依存関係ツリーを作成しています
状態情報を読み取っています... 完了
以下の追加パッケージがインストールされます:
  libpython-stdlib libpython2.7-minimal libpython2.7-stdlib mime-support python python-minimal
  python2.7 python2.7-minimal
提案パッケージ:
  python-doc python-tk python2.7-doc binutils binfmt-support
推奨パッケージ:
  file
以下のパッケージは「削除」されます:
  cron
以下のパッケージが新たにインストールされます:
  libpython-stdlib libpython2.7-minimal libpython2.7-stdlib mime-support python python-minimal
  python2.7 python2.7-minimal systemd-cron
...
\end{commandline}
\noindent
現状では /etc/cron.{hourly|daily|weekly|monthly} などをそれぞれ
\begin{commandline}
  % systemctl list-unit-files | grep cron
cron-update.path                       static
cron-daily.service                     static
cron-hourly.service                    static
cron-monthly.service                   static
cron-update.service                    static
cron-weekly.service                    static
cron.service                           masked
cron-daily.target                      static
cron-hourly.target                     static
cron-monthly.target                    static
cron-weekly.target                     static
cron.target                            enabled
cron-daily.timer                       static
cron-hourly.timer                      static
cron-monthly.timer                     static
cron-weekly.timer                      static
\end{commandline}
\noindent
といった塩梅に service と timer の units ファイルへ実行します。
中身は
\begin{commandline}
  % cat /lib/systemd/system/cron-hourly.timer
[Unit]
Description=systemd-cron hourly timer
Documentation=man:systemd.cron(7)
PartOf=cron.target
RefuseManualStart=yes
RefuseManualStop=yes

[Timer]
OnCalendar=hourly
Unit=cron-hourly.target
  % cat /lib/systemd/system/cron-hourly.service
[Unit]
Description=systemd-cron hourly script service
Documentation=man:systemd.cron(7)
PartOf=cron-hourly.target
RefuseManualStart=yes
RefuseManualStop=yes
ConditionDirectoryNotEmpty=/etc/cron.hourly

[Service]
Type=oneshot
ExecStart=/bin/run-parts --report /etc/cron.hourly
\end{commandline}
\noindent
となっており、まあ /etc/crontab に設定されている内容であれば、
ちゃんと実行できるのではないかなぁ、と期待します(が、どうなのかなぁ...)。

ちなみに「ユーザ毎の crontab がちゃんと migrate されているのか」も試して
みましたが、これは cron-update.{path,service} によっ
て /var/spool/cron/crontabls 以下のスクリプトが実行されている様で、
一応使えている様です。

ですが \verb|crontab -e| を叩いてみると
\begin{commandline}
  % crontab -e
  Traceback (most recent call last):
  File "/usr/bin/crontab", line 94, in <module>
    action(cron_file, args)
  File "/usr/bin/crontab", line 64, in edit
    with open(cron_file, 'r') as inp:
IOError: [Errno 2] No such file or directory: '/var/spool/cron/crontabs/root'
\end{commandline}
となるので、更新ができません。
この辺を上手く切り分けないと流石に移行は厳しそうですね。

\subsection{まとめ$\dots$?}

というわけで、systemd-{timesyncd,resolved,networkd,cron} について、
ごく簡単に触った感じをまとめてみました。
%
ちまたに溢れる systemd の解説では「もっとイロイロできるよ」と書いてあったりしますが、
version によって設定ファイルの書き方が変わっていたり、機能が追加/削除されていたりするので、
「頼れるのは man だけ」という、なかなか楽しい試行錯誤でした%
(逆に man は非常きちんとしているのが他のプログラムとは違うところですね)。

\renewcommand{\refname}{\normalsize{\textbf{参考文献}}}
\begin{thebibliography}{99}
\bibitem{佐々木(2014)}%
  佐々木洋平, 2014: Debian での systemd とのつきあい方,
  第85回 関西Debian勉強会, 2014年6月22日
\bibitem{systemd-timesyncd}
  \texttt{systemd-timesyncd(8), systemd-resolved(8), systemd-networkd(8)}
\end{thebibliography}


\dancersection{今後の予定}{Debian JP}

\subsection{関西Debian勉強会}
次回、第109回関西Debian勉強会は4月23日(土)か24(日)に開催予定です。
いつもの第4日曜日ではなく土曜日の開催を検討しております。

\subsection{東京エリアDebian勉強会}
次回、第139回東京エリアDebian勉強会は4月16日(土)に開催予定です。

%
% 冊子にするために、4の倍数にする必要がある。
% そのための調整
\dancersection{メモ}{}
\mbox{}\newpage
\mbox{}\newpage
\mbox{}\newpage

\pagebreak

\begin{center}
本資料のライセンスについて
\end{center}

本資料はフリー・ソフトウェアです。あなたは、Free Software
Foundation が公表したGNU GENERAL PUBLIC LICENSEの "バージョン２"もしくはそれ以降
が定める条項に従って本プログラムを再頒布または変更することができ
ます。

本プログラムは有用とは思いますが、頒布にあたっては、市場性及び特
定目的適合性についての暗黙の保証を含めて、いかなる保証も行ないま
せん。詳細についてはGNU GENERAL PUBLIC LICENSE をお読みください。

\begin{multicols}{2}
 \begin{fontsize}{6}{6}
 \begin{verbatim}
            GNU GENERAL PUBLIC LICENSE
               Version 2, June 1991

 Copyright (C) 1989, 1991 Free Software Foundation, Inc.
    51 Franklin St, Fifth Floor, Boston, MA  02110-1301  USA
 Everyone is permitted to copy and distribute verbatim copies
 of this license document, but changing it is not allowed.

                Preamble

  The licenses for most software are designed to take away your
freedom to share and change it.  By contrast, the GNU General Public
License is intended to guarantee your freedom to share and change free
software--to make sure the software is free for all its users.  This
General Public License applies to most of the Free Software
Foundation's software and to any other program whose authors commit to
using it.  (Some other Free Software Foundation software is covered by
the GNU Library General Public License instead.)  You can apply it to
your programs, too.

  When we speak of free software, we are referring to freedom, not
price.  Our General Public Licenses are designed to make sure that you
have the freedom to distribute copies of free software (and charge for
this service if you wish), that you receive source code or can get it
if you want it, that you can change the software or use pieces of it
in new free programs; and that you know you can do these things.

  To protect your rights, we need to make restrictions that forbid
anyone to deny you these rights or to ask you to surrender the rights.
These restrictions translate to certain responsibilities for you if you
distribute copies of the software, or if you modify it.

  For example, if you distribute copies of such a program, whether
gratis or for a fee, you must give the recipients all the rights that
you have.  You must make sure that they, too, receive or can get the
source code.  And you must show them these terms so they know their
rights.

  We protect your rights with two steps: (1) copyright the software, and
(2) offer you this license which gives you legal permission to copy,
distribute and/or modify the software.

  Also, for each author's protection and ours, we want to make certain
that everyone understands that there is no warranty for this free
software.  If the software is modified by someone else and passed on, we
want its recipients to know that what they have is not the original, so
that any problems introduced by others will not reflect on the original
authors' reputations.

  Finally, any free program is threatened constantly by software
patents.  We wish to avoid the danger that redistributors of a free
program will individually obtain patent licenses, in effect making the
program proprietary.  To prevent this, we have made it clear that any
patent must be licensed for everyone's free use or not licensed at all.

  The precise terms and conditions for copying, distribution and
modification follow.

            GNU GENERAL PUBLIC LICENSE
   TERMS AND CONDITIONS FOR COPYING, DISTRIBUTION AND MODIFICATION

  0. This License applies to any program or other work which contains
a notice placed by the copyright holder saying it may be distributed
under the terms of this General Public License.  The "Program", below,
refers to any such program or work, and a "work based on the Program"
means either the Program or any derivative work under copyright law:
that is to say, a work containing the Program or a portion of it,
either verbatim or with modifications and/or translated into another
language.  (Hereinafter, translation is included without limitation in
the term "modification".)  Each licensee is addressed as "you".

Activities other than copying, distribution and modification are not
covered by this License; they are outside its scope.  The act of
running the Program is not restricted, and the output from the Program
is covered only if its contents constitute a work based on the
Program (independent of having been made by running the Program).
Whether that is true depends on what the Program does.

  1. You may copy and distribute verbatim copies of the Program's
source code as you receive it, in any medium, provided that you
conspicuously and appropriately publish on each copy an appropriate
copyright notice and disclaimer of warranty; keep intact all the
notices that refer to this License and to the absence of any warranty;
and give any other recipients of the Program a copy of this License
along with the Program.

You may charge a fee for the physical act of transferring a copy, and
you may at your option offer warranty protection in exchange for a fee.

  2. You may modify your copy or copies of the Program or any portion
of it, thus forming a work based on the Program, and copy and
distribute such modifications or work under the terms of Section 1
above, provided that you also meet all of these conditions:

    a) You must cause the modified files to carry prominent notices
    stating that you changed the files and the date of any change.

    b) You must cause any work that you distribute or publish, that in
    whole or in part contains or is derived from the Program or any
    part thereof, to be licensed as a whole at no charge to all third
    parties under the terms of this License.

    c) If the modified program normally reads commands interactively
    when run, you must cause it, when started running for such
    interactive use in the most ordinary way, to print or display an
    announcement including an appropriate copyright notice and a
    notice that there is no warranty (or else, saying that you provide
    a warranty) and that users may redistribute the program under
    these conditions, and telling the user how to view a copy of this
    License.  (Exception: if the Program itself is interactive but
    does not normally print such an announcement, your work based on
    the Program is not required to print an announcement.)

These requirements apply to the modified work as a whole.  If
identifiable sections of that work are not derived from the Program,
and can be reasonably considered independent and separate works in
themselves, then this License, and its terms, do not apply to those
sections when you distribute them as separate works.  But when you
distribute the same sections as part of a whole which is a work based
on the Program, the distribution of the whole must be on the terms of
this License, whose permissions for other licensees extend to the
entire whole, and thus to each and every part regardless of who wrote it.

Thus, it is not the intent of this section to claim rights or contest
your rights to work written entirely by you; rather, the intent is to
exercise the right to control the distribution of derivative or
collective works based on the Program.

In addition, mere aggregation of another work not based on the Program
with the Program (or with a work based on the Program) on a volume of
a storage or distribution medium does not bring the other work under
the scope of this License.

  3. You may copy and distribute the Program (or a work based on it,
under Section 2) in object code or executable form under the terms of
Sections 1 and 2 above provided that you also do one of the following:

    a) Accompany it with the complete corresponding machine-readable
    source code, which must be distributed under the terms of Sections
    1 and 2 above on a medium customarily used for software interchange; or,

    b) Accompany it with a written offer, valid for at least three
    years, to give any third party, for a charge no more than your
    cost of physically performing source distribution, a complete
    machine-readable copy of the corresponding source code, to be
    distributed under the terms of Sections 1 and 2 above on a medium
    customarily used for software interchange; or,

    c) Accompany it with the information you received as to the offer
    to distribute corresponding source code.  (This alternative is
    allowed only for noncommercial distribution and only if you
    received the program in object code or executable form with such
    an offer, in accord with Subsection b above.)

The source code for a work means the preferred form of the work for
making modifications to it.  For an executable work, complete source
code means all the source code for all modules it contains, plus any
associated interface definition files, plus the scripts used to
control compilation and installation of the executable.  However, as a
special exception, the source code distributed need not include
anything that is normally distributed (in either source or binary
form) with the major components (compiler, kernel, and so on) of the
operating system on which the executable runs, unless that component
itself accompanies the executable.

If distribution of executable or object code is made by offering
access to copy from a designated place, then offering equivalent
access to copy the source code from the same place counts as
distribution of the source code, even though third parties are not
compelled to copy the source along with the object code.

  4. You may not copy, modify, sublicense, or distribute the Program
except as expressly provided under this License.  Any attempt
otherwise to copy, modify, sublicense or distribute the Program is
void, and will automatically terminate your rights under this License.
However, parties who have received copies, or rights, from you under
this License will not have their licenses terminated so long as such
parties remain in full compliance.

  5. You are not required to accept this License, since you have not
signed it.  However, nothing else grants you permission to modify or
distribute the Program or its derivative works.  These actions are
prohibited by law if you do not accept this License.  Therefore, by
modifying or distributing the Program (or any work based on the
Program), you indicate your acceptance of this License to do so, and
all its terms and conditions for copying, distributing or modifying
the Program or works based on it.

  6. Each time you redistribute the Program (or any work based on the
Program), the recipient automatically receives a license from the
original licensor to copy, distribute or modify the Program subject to
these terms and conditions.  You may not impose any further
restrictions on the recipients' exercise of the rights granted herein.
You are not responsible for enforcing compliance by third parties to
this License.

  7. If, as a consequence of a court judgment or allegation of patent
infringement or for any other reason (not limited to patent issues),
conditions are imposed on you (whether by court order, agreement or
otherwise) that contradict the conditions of this License, they do not
excuse you from the conditions of this License.  If you cannot
distribute so as to satisfy simultaneously your obligations under this
License and any other pertinent obligations, then as a consequence you
may not distribute the Program at all.  For example, if a patent
license would not permit royalty-free redistribution of the Program by
all those who receive copies directly or indirectly through you, then
the only way you could satisfy both it and this License would be to
refrain entirely from distribution of the Program.

If any portion of this section is held invalid or unenforceable under
any particular circumstance, the balance of the section is intended to
apply and the section as a whole is intended to apply in other
circumstances.

It is not the purpose of this section to induce you to infringe any
patents or other property right claims or to contest validity of any
such claims; this section has the sole purpose of protecting the
integrity of the free software distribution system, which is
implemented by public license practices.  Many people have made
generous contributions to the wide range of software distributed
through that system in reliance on consistent application of that
system; it is up to the author/donor to decide if he or she is willing
to distribute software through any other system and a licensee cannot
impose that choice.

This section is intended to make thoroughly clear what is believed to
be a consequence of the rest of this License.

  8. If the distribution and/or use of the Program is restricted in
certain countries either by patents or by copyrighted interfaces, the
original copyright holder who places the Program under this License
may add an explicit geographical distribution limitation excluding
those countries, so that distribution is permitted only in or among
countries not thus excluded.  In such case, this License incorporates
the limitation as if written in the body of this License.

  9. The Free Software Foundation may publish revised and/or new versions
of the General Public License from time to time.  Such new versions will
be similar in spirit to the present version, but may differ in detail to
address new problems or concerns.

Each version is given a distinguishing version number.  If the Program
specifies a version number of this License which applies to it and "any
later version", you have the option of following the terms and conditions
either of that version or of any later version published by the Free
Software Foundation.  If the Program does not specify a version number of
this License, you may choose any version ever published by the Free Software
Foundation.

  10. If you wish to incorporate parts of the Program into other free
programs whose distribution conditions are different, write to the author
to ask for permission.  For software which is copyrighted by the Free
Software Foundation, write to the Free Software Foundation; we sometimes
make exceptions for this.  Our decision will be guided by the two goals
of preserving the free status of all derivatives of our free software and
of promoting the sharing and reuse of software generally.

                NO WARRANTY

  11. BECAUSE THE PROGRAM IS LICENSED FREE OF CHARGE, THERE IS NO WARRANTY
FOR THE PROGRAM, TO THE EXTENT PERMITTED BY APPLICABLE LAW.  EXCEPT WHEN
OTHERWISE STATED IN WRITING THE COPYRIGHT HOLDERS AND/OR OTHER PARTIES
PROVIDE THE PROGRAM "AS IS" WITHOUT WARRANTY OF ANY KIND, EITHER EXPRESSED
OR IMPLIED, INCLUDING, BUT NOT LIMITED TO, THE IMPLIED WARRANTIES OF
MERCHANTABILITY AND FITNESS FOR A PARTICULAR PURPOSE.  THE ENTIRE RISK AS
TO THE QUALITY AND PERFORMANCE OF THE PROGRAM IS WITH YOU.  SHOULD THE
PROGRAM PROVE DEFECTIVE, YOU ASSUME THE COST OF ALL NECESSARY SERVICING,
REPAIR OR CORRECTION.

  12. IN NO EVENT UNLESS REQUIRED BY APPLICABLE LAW OR AGREED TO IN WRITING
WILL ANY COPYRIGHT HOLDER, OR ANY OTHER PARTY WHO MAY MODIFY AND/OR
REDISTRIBUTE THE PROGRAM AS PERMITTED ABOVE, BE LIABLE TO YOU FOR DAMAGES,
INCLUDING ANY GENERAL, SPECIAL, INCIDENTAL OR CONSEQUENTIAL DAMAGES ARISING
OUT OF THE USE OR INABILITY TO USE THE PROGRAM (INCLUDING BUT NOT LIMITED
TO LOSS OF DATA OR DATA BEING RENDERED INACCURATE OR LOSSES SUSTAINED BY
YOU OR THIRD PARTIES OR A FAILURE OF THE PROGRAM TO OPERATE WITH ANY OTHER
PROGRAMS), EVEN IF SUCH HOLDER OR OTHER PARTY HAS BEEN ADVISED OF THE
POSSIBILITY OF SUCH DAMAGES.

             END OF TERMS AND CONDITIONS

        How to Apply These Terms to Your New Programs

  If you develop a new program, and you want it to be of the greatest
possible use to the public, the best way to achieve this is to make it
free software which everyone can redistribute and change under these terms.

  To do so, attach the following notices to the program.  It is safest
to attach them to the start of each source file to most effectively
convey the exclusion of warranty; and each file should have at least
the "copyright" line and a pointer to where the full notice is found.

    <one line to give the program's name and a brief idea of what it does.>
    Copyright (C) <year>  <name of author>

    This program is free software; you can redistribute it and/or modify
    it under the terms of the GNU General Public License as published by
    the Free Software Foundation; either version 2 of the License, or
    (at your option) any later version.

    This program is distributed in the hope that it will be useful,
    but WITHOUT ANY WARRANTY; without even the implied warranty of
    MERCHANTABILITY or FITNESS FOR A PARTICULAR PURPOSE.  See the
    GNU General Public License for more details.

    You should have received a copy of the GNU General Public License
    along with this program; if not, write to the Free Software
    Foundation, Inc., 51 Franklin St, Fifth Floor, Boston, MA  02110-1301 USA


Also add information on how to contact you by electronic and paper mail.

If the program is interactive, make it output a short notice like this
when it starts in an interactive mode:

    Gnomovision version 69, Copyright (C) year  name of author
    Gnomovision comes with ABSOLUTELY NO WARRANTY; for details type `show w'.
    This is free software, and you are welcome to redistribute it
    under certain conditions; type `show c' for details.

The hypothetical commands `show w' and `show c' should show the appropriate
parts of the General Public License.  Of course, the commands you use may
be called something other than `show w' and `show c'; they could even be
mouse-clicks or menu items--whatever suits your program.

You should also get your employer (if you work as a programmer) or your
school, if any, to sign a "copyright disclaimer" for the program, if
necessary.  Here is a sample; alter the names:

  Yoyodyne, Inc., hereby disclaims all copyright interest in the program
  `Gnomovision' (which makes passes at compilers) written by James Hacker.

  <signature of Ty Coon>, 1 April 1989
  Ty Coon, President of Vice

This General Public License does not permit incorporating your program into
proprietary programs.  If your program is a subroutine library, you may
consider it more useful to permit linking proprietary applications with the
library.  If this is what you want to do, use the GNU Library General
Public License instead of this License.
 \end{verbatim}
 \end{fontsize}
\end{multicols}

\begin{center}
ソースコードについて
\end{center}

このプログラムは tex で記述されたものです。ソースコードは
\begin{center}
  \url{git://anonscm.debian.org/tokyodebian/monthly-report.git}
\end{center}
から取得できます。

\begin{center}
Debian オープンユーズロゴ ライセンス
\end{center}

\begin{multicols}{2}
 \begin{fontsize}{6}{6}
 \begin{verbatim}

Copyright (c) 1999 Software in the Public Interest
Permission is hereby granted, free of charge, to any person
obtaining a copy of this software and associated documentation
files (the "Software"), to deal in the Software without restriction,
including without limitation the rights to use, copy, modify, merge,
publish, distribute, sublicense, and/or sell copies of the Software,
and to permit persons to whom the Software is furnished to do so,
subject to the following conditions:

The above copyright notice and this permission notice shall be
included in all copies or substantial portions of the Software.

THE SOFTWARE IS PROVIDED "AS IS", WITHOUT WARRANTY OF ANY
KIND, EXPRESS OR IMPLIED, INCLUDING BUT NOT LIMITED TO THE
WARRANTIES OF MERCHANTABILITY, FITNESS FOR A PARTICULAR PURPOSE AND
NONINFRINGEMENT. IN NO EVENT SHALL THE AUTHORS OR COPYRIGHT HOLDERS
BE LIABLE FOR ANY CLAIM, DAMAGES OR OTHER LIABILITY, WHETHER IN
AN ACTION OF CONTRACT, TORT OR OTHERWISE, ARISING FROM, OUT OF OR
IN CONNECTION WITH THE SOFTWARE OR THE USE OR OTHER DEALINGS IN
THE SOFTWARE.
 \end{verbatim}
 \end{fontsize}
\end{multicols}

\printindex
%\cleartooddpage

 \begin{minipage}[b]{0.2\hsize}
  \rotatebox{90}{\fontsize{80}{80} {\gt 関西 Debian 勉強会} }
 \end{minipage}
 \begin{minipage}[b]{0.8\hsize}

 \vspace*{15cm}
 \rule{\hsize}{1mm}
 \vspace{2mm}
 \includegraphics[width=2cm]{image200502/openlogo-nd.eps}
 \noindent \Large \bfseries{Debian 勉強会資料}\\ \\
 \noindent \normalfont \debmtgyear{}年\debmtgmonth{}月\debmtgdate{}日 \hspace{5mm}  初版第1刷発行\\
 \noindent \normalfont 関西 Debian 勉強会 （編集・印刷・発行）\\
 \rule{\hsize}{1mm}
 \end{minipage}

\end{document}
%%% Local Variables:
%%% skk-kutouten-type: jp
%%% End: